\section{Tiered Compute Runtime}

\label{sec:tiers}

\subsection{Motivation}

A key insight of the UAH architecture is that different agent tasks require
radically different compute environments. Answering a factual question requires
only an LLM inference call (milliseconds, cents). Resolving a GitHub issue
requires a sandboxed code execution environment (seconds, dollars). Automating
a complex multi-application workflow requires a full virtual desktop with GUI
access (minutes, tens of dollars). Building and deploying a production service
requires persistent cloud compute (hours, hundreds of dollars).

Traditional agent products hard-code a single compute tier. Devin always
provisions a full VM. Claude Code always runs in the user's terminal. The UAH
dynamically selects and upgrades compute tier based on task requirements,
with automatic context migration on tier transitions.

\subsection{The Five-Tier Hierarchy}

We define the Compute Tier as an element of the ordered set:
\begin{equation}
    \mathcal{T} = \left\{T_0 \prec T_1 \prec T_2 \prec T_3 \prec T_4\right\},
\end{equation}
where $T_i \prec T_j$ means $T_i$ is strictly less capable and less costly
than $T_j$.

\paragraph{Tier 0: Shared Gateway (Message Routing).}
No persistent execution environment. The agent receives a message, invokes
the LLM, calls stateless tools (web search, document retrieval, computation),
and returns a response. Latency: 1--5 seconds. Cost: $10^{-3}$--$10^{-2}$
USD per interaction.

\paragraph{Tier 1: Terminal Pod (Isolated Container).}
A dedicated Docker container with a filesystem, shell, and standard
Unix tools. Supports code execution, file I/O, and network requests within
a sandboxed namespace. Lifetime: up to 4 hours. Cost: $\sim 0.02$ USD/hour.
Used for: code generation and testing, data analysis scripts, automated
report generation.

\paragraph{Tier 2: Operative Desktop (VNC + GUI Environment).}
A full Linux desktop environment (Ubuntu 22.04) with VNC access. Provides
a graphical display server (Xvfb), browser automation (Playwright/Chromium),
and GUI applications. The Hanzo Operative framework~\cite{hanzo2024operative}
manages screen understanding and action execution. Cost: $\sim 0.10$ USD/hour.
Used for: web automation, GUI testing, multi-application workflows.

\paragraph{Tier 3: Cloud VM (Full Compute).}
A dedicated cloud virtual machine (4--64 vCPUs, 8--256 GB RAM) with
persistent storage, public IP, and full system access. Supports long-running
processes, custom software installation, and external network services. Cost:
$0.20$--$5.00$ USD/hour. Used for: ML training runs, system administration,
multi-service deployments.

\paragraph{Tier 4: Local Desktop (Hardware Access).}
The user's physical desktop or laptop machine, registered with the gateway
via a local agent daemon. Provides direct hardware access (GPU, USB devices,
camera, microphone) with no virtualization overhead. Registered machines
communicate with the cloud gateway via an encrypted tunnel, allowing cloud
orchestration of local compute. Cost: user's own hardware (compute is free,
tunnel bandwidth $\sim 0.01$ USD/GB).

\subsection{Tier Transition State Machine}

We model tier management as a state machine $(Q, \Sigma, \delta, q_0, F)$
where:
\begin{itemize}[leftmargin=1.5em]
    \item $Q = \{q_0, q_1, q_2, q_3, q_4\}$ are the tier states.
    \item $\Sigma = \{\texttt{upgrade}, \texttt{downgrade}, \texttt{migrate}\}$
          are transitions.
    \item $\delta : Q \times \Sigma \to Q$ is the transition function.
    \item $q_0 = T_0$ is the initial state (all sessions start at Tier 0).
    \item $F = Q$ (all tiers are accepting; the machine is reset on session end).
\end{itemize}

Each transition carries a \textbf{precondition} and \textbf{migration protocol}:

\begin{definition}[Tier Upgrade]
The transition $\delta(T_i, \texttt{upgrade}) = T_j$ ($j > i$) is valid iff:
\begin{enumerate}[leftmargin=1.5em, label=(\roman*)]
    \item $\text{UserConsentLevel}(T_j) \leq \text{UserGrantedConsent}()$
    \item $\text{Budget}(\text{session}) \geq \text{EstimatedCost}(T_j, \text{task})$
    \item $\text{TaskComplexityScore}(\text{task}) \geq \theta_j$
\end{enumerate}
where $\theta_j$ is the complexity threshold for tier $T_j$.
\end{definition}

\begin{definition}[Context Migration Protocol]
On upgrade from $T_i$ to $T_j$, the context migration function
$\mu_{i \to j} : \text{Context}_{T_i} \to \text{Context}_{T_j}$ performs:
\begin{enumerate}[leftmargin=1.5em, label=(\roman*)]
    \item Serialize current working state (files, variables, conversation).
    \item Provision the target compute environment $T_j$.
    \item Transfer state via the Hanzo context bundle format (CBOR-encoded,
          AES-256-GCM encrypted).
    \item Resume execution with injected system prompt summarizing the prior
          tier context.
    \item Verify environment readiness with a liveness probe before returning
          control to the agent.
\end{enumerate}
\end{definition}

The migration is atomic from the agent's perspective: the agent issues a
single tool call (\texttt{upgrade\_tier}) and receives an updated
\texttt{ExecutionEnvironment} handle. The migration latency is bounded:
Tier 1 provisioning completes in $\leq 8$ seconds (pre-warmed pool),
Tier 2 in $\leq 30$ seconds, Tier 3 in $\leq 90$ seconds.

\begin{algorithm}[H]
\caption{Adaptive Tier Selection}
\label{alg:tier}
\begin{algorithmic}[1]
\Require Task description $\tau$, session context $\mathcal{C}$, policy $\pi$
\Ensure Execution environment $e \in \{T_0, \ldots, T_4\}$
\State $s_\tau \gets \text{ComplexityScore}(\tau, \mathcal{C})$
\Comment{LLM-based scoring}
\State $t \gets \min\{j : \theta_j \leq s_\tau\}$
\Comment{Minimum sufficient tier}
\State $b \gets \text{SessionBudget}(\mathcal{C})$
\If{$\text{EstimatedCost}(T_t, \tau) > b$}
    \State $t \gets \max\{j : \text{EstimatedCost}(T_j, \tau) \leq b\}$
    \Comment{Budget-constrained fallback}
\EndIf
\State $\text{consent} \gets \text{VerifyConsent}(T_t, \mathcal{C}.\text{user})$
\If{$\neg\,\text{consent}$}
    \State $\textbf{return}\ \text{RequestConsentAndRetry}(T_t, \tau)$
\EndIf
\State $e \gets \text{ProvisionTier}(T_t, \mathcal{C})$
\State $\text{Emit}(\text{TIER\_SELECTED}, \{t, s_\tau, b, \mathcal{C}.\text{session\_id}\})$
\State \Return $e$
\end{algorithmic}
\end{algorithm}

\subsection{Downgrade and Session Cleanup}

When task complexity decreases (a common occurrence in long sessions), the
runtime can downgrade the tier to reduce cost. A downgrade $\delta(T_i,
\texttt{downgrade}) = T_j$ ($j < i$) is triggered when:
\begin{equation}
    \text{IdleTime}(T_i) \geq \Delta_{\text{idle}} \quad \text{and} \quad
    \text{PendingWork}(\text{session}) \subseteq \text{Capabilities}(T_j).
\end{equation}
The downgrade uses the same migration protocol, with the added step of
flushing all persistent state to external storage (MinIO object store) before
releasing the compute environment.

% ============================================================================
% 5. CROSS-LAYER OPTIMIZATION
% ============================================================================
